\documentclass[11pt]{article}
\usepackage{amsmath,amssymb,amsthm}
\usepackage[margin=1in]{geometry}
\usepackage{url}
\urlstyle{tt}

\newtheorem{theorem}{Theorem}
\newtheorem{lemma}[theorem]{Lemma}
\newtheorem{corollary}[theorem]{Corollary}
\newtheorem{proposition}[theorem]{Proposition}
\theoremstyle{remark}
\newtheorem*{remark}{Remark}

\title{Dense cyclotomic subgroups of the Lorentz group\\
and the Collins--Perez--Sudarsky obstruction}
\author{A.~M.~Thorn}
\date{2026}

\begin{document}
\maketitle

\begin{abstract}
We construct a finitely generated subgroup
$\Gamma = \langle A_5, B_J\rangle$ of $\mathrm{SO}^+(3,1)^0$ from the
rotational icosahedral group $A_5$ and a single Lorentz boost $B_J$
of rapidity $\log\varphi$, where $\varphi = (1+\sqrt{5})/2$ arises
from the cyclotomic unit $J = 1 + \zeta_5^2 \in \mathbb{Z}[\zeta_5]$.
Density of $\Gamma$ in the identity component is proved in two steps.
First, an explicit commutator of two conjugate boosts is shown to be
a pure elliptic Lorentz element with rotation cosine
$c_\omega = (31 + 8\sqrt{5})/50$, satisfying the irreducible quadratic
$2500\,x^2 - 3100\,x + 641$ over $\mathbb{Q}$; Niven's theorem and a
direct numerical comparison rule out rational $\omega/\pi$, so the
cyclic group it generates is dense in a one-parameter elliptic
subgroup. Second, the Lie algebra of the identity component of the
closure is $\mathrm{Ad}(\Gamma)$-invariant; combining the
irreducibility of the $A_5$-action on
$\mathfrak{so}(3,1) = \mathfrak{k} \oplus \mathfrak{p}$ with the
bracket structure forces this Lie algebra to be all of
$\mathfrak{so}(3,1)$. Topologically, every continuous
$\Gamma$-invariant functional on a configuration space with
continuous $\mathrm{SO}^+(3,1)^0$-action is automatically Lorentz
invariant. Under standard regularity assumptions on the dynamics,
measure, regulator, and coarse-graining, the resulting Wilson
effective action of any substrate field theory with $\Gamma$ in its
symmetry group cannot contain Lorentz-violating operators, addressing
the Collins--Perez--Sudarsky--Urrutia--Vucetich (2004) radiative
fine-tuning obstruction at its premise.
\end{abstract}


\section{Introduction}

Discrete-substrate models of spacetime face a structural obstruction
identified by Collins, Perez, Sudarsky, Urrutia and Vucetich
\cite{CPSUV04}. A discrete substrate breaks Lorentz invariance at the
ultraviolet cutoff $\Lambda \sim E_{\mathrm{Planck}}$, and radiative
corrections in the effective field theory below $\Lambda$ generically
inflate this ultraviolet violation into infrared coefficients of order
$\Lambda^2 / 16\pi^2$. To match the observed bounds on Lorentz
violation \cite{Mattingly05,LiberatiReview}, the ultraviolet
coefficients must be fine-tuned to a precision of $10^{-20}$ or finer,
with no symmetry known to enforce the tuning. Most discrete substrates
examined in the literature fall to this argument: cubic lattice gauge
theories, naive lattice formulations of loop quantum gravity, and
Lorentz-violating variants of doubly-special relativity. Causal-set
theory \cite{BHS87} escapes by appeal to Poisson sprinkling and a
statistical Lorentz invariance.

The mechanism of the Collins--Perez--Sudarsky obstruction is a mismatch
between the ultraviolet symmetry $G_{\mathrm{UV}}$ of the substrate
and the infrared symmetry $\mathrm{SO}^+(3,1)^0$ of the effective
theory. $G_{\mathrm{UV}}$-invariant operators that are not also
Lorentz-invariant are present in the effective Lagrangian, and
radiative corrections generate them with their natural coefficients.
Fine-tuning is required to suppress them.

This paper observes that the obstruction does not apply when
$G_{\mathrm{UV}}$ is a dense subgroup of $\mathrm{SO}^+(3,1)^0$. In
that case, the set of $G_{\mathrm{UV}}$-invariant continuous
functionals on a Lorentz-equivariant field space coincides with the
set of $\mathrm{SO}^+(3,1)^0$-invariant continuous functionals. No
Lorentz-violating operator can occur; there is nothing to tune.

We exhibit a concrete example. Let $K = \mathbb{Q}(\zeta_5)$ with
$\zeta_5 = e^{2\pi i/5}$, and let $J = 1 + \zeta_5^2 \in
\mathcal{O}_K$. Galois norm and embeddings give the four conjugate
moduli $(|\sigma_1 J|, |\sigma_2 J|, |\sigma_3 J|, |\sigma_4 J|) =
(\varphi^{-1}, \varphi, \varphi, \varphi^{-1})$ where $\varphi =
(1+\sqrt{5})/2$, so $J$ is a unit. We use $J$ to fix the rapidity
$\rho = \log\varphi$ of a Lorentz boost $B_J \in \mathrm{SO}^+(3,1)^0$
along a chosen axis: $\rho = \log\varphi$ is the positive
logarithmic absolute value occurring among the Galois conjugates of
$J$ (it occurs twice, for $k = 2$ and $k = 3$). Let $A_5 \subset \mathrm{SO}(3)$ be the
rotational icosahedral group (the image of the binary icosahedral
group $2I \subset \mathrm{SU}(2)$ under the spin double cover), embedded
in $\mathrm{SO}^+(3,1)^0$ via the block embedding $g \mapsto
\mathrm{diag}(1, g)$. The group $\Gamma = \langle A_5, B_J\rangle
\subset \mathrm{SO}^+(3,1)^0$ has infinite order. Our main theorem
establishes that its closure is the full identity component.

The proof has two steps. An explicit commutator of conjugate boosts
yields an infinite-order elliptic Lorentz element whose rotation
cosine $c_\omega$ has degree two over $\mathbb{Q}$; Niven's theorem
restricts a putative rational $\omega/\pi$ to a finite list of
candidate cosines, and direct comparison rules them all out. The
cyclic group generated by this element is therefore dense in a
one-parameter elliptic subgroup of $\mathrm{SO}^+(3,1)^0$. The Lie
algebra of the identity component of the closure is then
$\mathrm{Ad}(\Gamma)$-invariant, and the irreducibility of $A_5$ on
$\mathfrak{so}(3,1) = \mathfrak{k} \oplus \mathfrak{p}$ together with
the bracket structure forces this Lie algebra to be all of
$\mathfrak{so}(3,1)$.

The Lorentz protection corollary is a topological consequence of
density: under standard continuity assumptions on the effective
dynamics, every continuous $\Gamma$-invariant functional is
$\mathrm{SO}^+(3,1)^0$-invariant.

A separate computational appendix verifies an unrelated but
useful Molien-series fact: $A_5$-invariant spatial couplings on
$\mathbb{R}^3$ have first angular anisotropy at degree six, two orders
later than cubic-symmetric couplings. This is a consequence of the
classical icosahedral invariant theory of Klein \cite{Klein1884} and
serves as a numerical check of the spatial-symmetry features of the
substrate, not of Lorentz protection itself. All scripts are released
as a companion artifact \cite{ALPCode}.


\section{Construction}

Let $K = \mathbb{Q}(\zeta_5)$, $\mathcal{O}_K = \mathbb{Z}[\zeta_5]$,
and $\varphi = (1+\sqrt{5})/2$. The Galois group
$\mathrm{Gal}(K/\mathbb{Q})$ is cyclic of order four, generated by
$\sigma_2: \zeta_5 \mapsto \zeta_5^2$. Set
\[
  J \;=\; 1 + \zeta_5^2 \;\in\; \mathcal{O}_K.
\]
Direct computation gives the conjugate moduli
$(|\sigma_1 J|, |\sigma_2 J|, |\sigma_3 J|, |\sigma_4 J|) =
(\varphi^{-1}, \varphi, \varphi, \varphi^{-1})$ and norm
$N_{K/\mathbb{Q}}(J) = 1$, so $J$ is a unit. The four real numbers
$\log|\sigma_k J|$ form the multiset
\[
  \bigl\{\, -\log\varphi,\; +\log\varphi,\; +\log\varphi,\;
    -\log\varphi \,\bigr\},
\]
and the positive logarithmic absolute value occurring among the
Galois conjugates is $\log\varphi$, with multiplicity two. We use this real positive
unit-logarithm as the rapidity $\rho = \log\varphi$ of a Lorentz
boost $B_J \in \mathrm{SO}^+(3,1)^0$ along the chosen axis $\hat{z}$.

In coordinates,
\[
  B_J \;=\;
  \begin{pmatrix}
    \cosh\rho & 0 & 0 & \sinh\rho \\
    0 & 1 & 0 & 0 \\
    0 & 0 & 1 & 0 \\
    \sinh\rho & 0 & 0 & \cosh\rho
  \end{pmatrix}
  \;=\;
  \begin{pmatrix}
    \sqrt{5}/2 & 0 & 0 & 1/2 \\
    0 & 1 & 0 & 0 \\
    0 & 0 & 1 & 0 \\
    1/2 & 0 & 0 & \sqrt{5}/2
  \end{pmatrix}.
\]
Note that $\cosh\rho = \sqrt{5}/2$ and $\sinh\rho = 1/2$ both lie in
$\mathbb{Q}(\sqrt{5})$, so $B_J$ has entries in $\mathbb{Q}(\sqrt{5})$.

The rotational icosahedral group $I = A_5 \subset \mathrm{SO}(3)$ is
the simple non-abelian group of order $60$, the image of the binary
icosahedral group $2I \subset \mathrm{SU}(2)$ under the spin double
cover. We embed $A_5$ into $\mathrm{SO}^+(3,1)^0$ as the spatial
rotation block: $g \mapsto \mathrm{diag}(1, g)$. Define
\[
  \Gamma \;=\; \bigl\langle\, A_5,\, B_J \,\bigr\rangle \;\leq\;
    \mathrm{SO}^+(3,1)^0.
\]


\section{Density theorem}

\begin{theorem}\label{thm:density}
The topological closure $\overline{\Gamma}$ in $\mathrm{SO}^+(3,1)^0$
equals the whole identity component.
\end{theorem}

\begin{proof}
Let $H = \overline{\Gamma}$ and let $H^0$ be its identity component.
We show $\mathrm{Lie}(H^0) = \mathfrak{so}(3,1)$ in two steps.

\emph{Step 1: an infinite-order elliptic element.}
Identify $\mathbb{R}^3$ with the spatial slice of $\mathbb{R}^{1,3}$
and place an icosahedron inscribed in the unit sphere with vertices
forming a single $A_5$-orbit of twelve directions. The standard
choice
\[
  v_1 = (0, 1, \varphi), \qquad v_2 = (1, \varphi, 0)
\]
gives two adjacent vertices with $\|v_1\|^2 = \|v_2\|^2 = \varphi+2$
and $v_1 \cdot v_2 = \varphi$, so the angle $\alpha$ between them
satisfies
\[
  \cos\alpha \;=\; \frac{v_1 \cdot v_2}{\|v_1\|\,\|v_2\|} \;=\;
  \frac{\varphi}{\varphi+2} \;=\; \frac{1}{\sqrt{5}}.
\]
The order-three rotation
\[
  g \;=\; \begin{pmatrix}
    -1/2 & (1-\sqrt{5})/4 & (1+\sqrt{5})/4 \\
    (-1+\sqrt{5})/4 & (1+\sqrt{5})/4 & 1/2 \\
    -(1+\sqrt{5})/4 & 1/2 & (1-\sqrt{5})/4
  \end{pmatrix}
\]
satisfies $g v_1 = v_2$, $g^T g = I$, $\det g = 1$, and $g^3 = I$.
A direct check shows that $g$ permutes the set of twelve standard
icosahedron vertices $\{(0, \pm 1, \pm\varphi), (\pm 1, \pm\varphi, 0),
(\pm\varphi, 0, \pm 1)\}$, with orbit structure four disjoint
$3$-cycles. Since $g$ permutes the vertices of the standard
icosahedron and has $\det g = 1$, it lies in the rotational symmetry
group of that icosahedron, which is the rotational icosahedral group
$I \cong A_5$. Geometrically, $g$ is the order-three rotation about
the centre of the icosahedral face shared between $v_1$ and $v_2$.

Density of a subgroup is invariant under conjugation in
$\mathrm{SO}^+(3,1)^0$: if $R \in \mathrm{SO}^+(3,1)^0$, then
$R \Gamma R^{-1}$ is dense if and only if $\Gamma$ is. After a
spatial conjugation by $R \in \mathrm{SO}(3) \subset
\mathrm{SO}^+(3,1)^0$ sending $v_1 / \|v_1\|$ to the third coordinate
axis, we may therefore take the boost axis to be
$\hat{z} = v_1 / \|v_1\|$. The other vertex direction is then
$\hat{n} = v_2 / \|v_2\|$, and we set $B_g = g B_J g^{-1}$, the boost
of rapidity $\rho = \log\varphi$ along $\hat{n}$. The commutator
\[
  C \;=\; B_J B_g B_J^{-1} B_g^{-1} \;\in\; \Gamma.
\]
Direct calculation in $\mathbb{Q}(\sqrt{5})$ gives the trace
\[
  \mathrm{tr}\,C \;=\; 2 \;+\; 2 c_\omega, \qquad
  c_\omega \;=\; \frac{31 + 8\sqrt{5}}{50},
\]
and the characteristic polynomial
\[
  \chi_C(\lambda) \;=\; (\lambda - 1)^2 \bigl(\lambda^2 - 2 c_\omega
  \lambda + 1\bigr).
\]
The eigenvalues are $1, 1, e^{i\omega}, e^{-i\omega}$ with
$\cos\omega = c_\omega$.

The characteristic polynomial alone is not quite enough to classify
$C$: one must also exclude a non-trivial Jordan block at the
eigenvalue~$1$, which would make $C$ parabolic rather than elliptic.
The same exact computation gives
\[
  \mathrm{rank}(C - I) \;=\; 2, \qquad \dim\ker(C - I) \;=\; 2,
\]
and the restriction of the Lorentz form $\eta$ to the two-dimensional
fixed space $\ker(C - I)$ has Gram matrix
\[
  G \;=\; \begin{pmatrix}
    -1 & 0 \\ 0 & (13 + \sqrt{5})/2
  \end{pmatrix},
  \qquad
  \det G \;=\; -\frac{13 + \sqrt{5}}{2} \;<\; 0,
\]
so $\eta$ has signature $(1, 1)$ on the fixed plane. Therefore $C$
is conjugate in $\mathrm{SO}^+(3,1)^0$ to a spatial rotation by angle
$\omega$ about the spacelike axis perpendicular to its fixed plane.
The trace, characteristic polynomial, fixed-plane Gram matrix, rank
of $C - I$, and minimal polynomial of $c_\omega$ are all reproduced
from these definitions in symbolic and 200-digit precision arithmetic
by claims [1], [2], [9] of the companion artifact \cite{ALPCode}.

The cosine $c_\omega$ has minimal polynomial
\[
  2500 x^2 \;-\; 3100 x \;+\; 641 \;=\; 0
\]
over $\mathbb{Q}$, with discriminant
$\Delta = 3100^2 - 4 \cdot 2500 \cdot 641 = 3\,200\,000 = 800^2 \cdot 5$,
non-square in $\mathbb{Q}$. The polynomial is irreducible,
$[\mathbb{Q}(c_\omega) : \mathbb{Q}] = 2$, and $c_\omega \in
\mathbb{Q}(\sqrt{5})$.

\begin{lemma}[Cyclotomic cosine exclusion]\label{lem:cce}
Let $c \in \mathbb{R}$ have $[\mathbb{Q}(c) : \mathbb{Q}] = 2$.
Suppose $c = \cos\omega$ with $\omega / \pi \in \mathbb{Q}$. Then
$\zeta := e^{i\omega}$ is a root of unity of some order $m \geq 1$,
and $2c = \zeta + \zeta^{-1}$ lies in the maximal real subfield of
$\mathbb{Q}(\zeta_m)$, which has degree $\phi(m)/2$ over
$\mathbb{Q}$. The constraint $\phi(m)/2 = 2$ forces
\[
  m \;\in\; \{5,\,8,\,10,\,12\}.
\]
The corresponding cosine values are
\[
  m = 5: \;\;\pm\tfrac{\sqrt{5}-1}{4},\; \mp\tfrac{1+\sqrt{5}}{4}; \qquad
  m = 8: \;\;\pm\tfrac{\sqrt{2}}{2}; \\
  m = 10: \;\;\pm\tfrac{1+\sqrt{5}}{4},\; \mp\tfrac{\sqrt{5}-1}{4}; \qquad
  m = 12: \;\;\pm\tfrac{\sqrt{3}}{2},
\]
giving the eight distinct candidate values
\[
  c \in \Bigl\{\pm\tfrac{\sqrt{2}}{2},\; \pm\tfrac{\sqrt{3}}{2},\;
    \pm\tfrac{1+\sqrt{5}}{4},\; \pm\tfrac{\sqrt{5}-1}{4}\Bigr\}.
\]
The lower cyclotomic orders $m \in \{1, 2, 3, 4, 6\}$ give cosine
values $\{\pm 1, 0, \pm 1/2\}$ in $\mathbb{Q}$ and are excluded by
$[\mathbb{Q}(c) : \mathbb{Q}] = 2$.
\end{lemma}

This is a standard consequence of cyclotomic theory; the rational
multiples-of-$\pi$ case also follows from Niven's theorem
\cite{Niven}.

Applying Lemma~\ref{lem:cce} to $c_\omega = (31 + 8\sqrt{5})/50$,
with degree two over $\mathbb{Q}$ established above: the values
$\pm\sqrt{2}/2, \pm\sqrt{3}/2$ do not lie in $\mathbb{Q}(\sqrt{5})$
(since $\sqrt{2}, \sqrt{3} \notin \mathbb{Q}(\sqrt{5})$). The four
remaining values $\pm(1+\sqrt{5})/4 \approx \pm 0.8090$ and
$\pm(\sqrt{5}-1)/4 \approx \pm 0.3090$ are numerically distinct from
$c_\omega \approx 0.9778$. Hence $\omega/\pi \notin \mathbb{Q}$.

The cyclic group $\langle C\rangle$ is therefore infinite, and its
topological closure in $\mathrm{SO}^+(3,1)^0$ is the unique
one-parameter elliptic subgroup $T_C \cong S^1$ of rotations sharing
$C$'s fixed Lorentz frame. We have $T_C \subseteq H^0$, so
$\dim H^0 \geq 1$, in particular $\mathfrak{h} = \mathrm{Lie}(H^0) \neq
\{0\}$.

\emph{Step 2: Lie algebra normality.}
Since $H^0$ is normal in $H$ and $\Gamma \subseteq H$, the Lie
algebra $\mathfrak{h}$ is invariant under $\mathrm{Ad}(\Gamma)$, in
particular under $\mathrm{Ad}(A_5)$ and $\mathrm{Ad}(B_J)$. Decompose
\[
  \mathfrak{so}(3,1) \;=\; \mathfrak{k} \;\oplus\; \mathfrak{p},
\]
with $\mathfrak{k} \cong \mathfrak{so}(3)$ generated by rotation
generators $J(v)$ and $\mathfrak{p}$ generated by boost generators
$K(v)$, for $v \in \mathbb{R}^3$. Both $\mathfrak{k}$ and
$\mathfrak{p}$ carry the standard three-dimensional irreducible
$A_5$-representation $\rho_3$, so as $A_5$-modules
\[
  \mathfrak{so}(3,1) \;\cong\; \rho_3 \oplus \rho_3.
\]
The real centralizer of the $A_5$-action on $\mathbb{R}^3$ is
$\mathbb{R} \cdot I$: a linear endomorphism commuting with all
icosahedral rotations must preserve each fivefold rotation axis, the
six fivefold axes span $\mathbb{R}^3$ in pairs, and matching scalars
across axes is forced by the action of any threefold rotation
permuting them. Hence $\rho_3$ is absolutely irreducible over
$\mathbb{R}$, of real type rather than complex or quaternionic.

By Schur's lemma applied over $\mathbb{R}$ to $\rho_3 \oplus \rho_3$,
every nonzero $A_5$-invariant proper $\mathbb{R}$-subspace of
$\mathfrak{so}(3,1)$ has dimension three and the form
\[
  L_\lambda \;=\; \bigl\{\, J(v) + \lambda\, K(v) : v \in
  \mathbb{R}^3 \,\bigr\}, \qquad \lambda \in \mathbb{R} \cup
  \{\infty\},
\]
where $L_\infty := \mathfrak{p}$. There are no $A_5$-invariant
subspaces of dimension one, two, four, or five, since $\rho_3$ is
absolutely irreducible and $\mathfrak{so}(3,1)$ contains it with
multiplicity exactly two.

The Lie brackets in $\mathfrak{so}(3,1)$ are
\[
  [J(u), J(v)] = J(u\times v), \quad
  [J(u), K(v)] = K(u\times v), \quad
  [K(u), K(v)] = -J(u\times v).
\]
For $X = J(u) + \lambda K(u)$ and $Y = J(v) + \lambda K(v)$,
\[
  [X, Y] \;=\; (1 - \lambda^2)\, J(u\times v) \;+\; 2\lambda\,
  K(u\times v).
\]
Requiring $[X, Y] \in L_\lambda$ means there exists $w \in
\mathbb{R}^3$ with $J(w) = (1-\lambda^2) J(u\times v)$ and $\lambda
K(w) = 2\lambda K(u\times v)$. The first gives $w = (1-\lambda^2)
(u \times v)$, and substituting into the second demands $\lambda(1 -
\lambda^2)(u\times v) = 2\lambda (u\times v)$ for all $u, v$ with
$u\times v \neq 0$. This yields $\lambda(1+\lambda^2) = 0$, whose
only real solution is $\lambda = 0$. Hence $L_0 = \mathfrak{k}$ is
the unique three-dimensional $A_5$-invariant Lie subalgebra. The
remaining case $L_\infty = \mathfrak{p}$ is not closed under brackets
either, since $[K(u), K(v)] = -J(u\times v) \in \mathfrak{k}$ lies
outside $\mathfrak{p}$. Therefore the only $A_5$-invariant Lie
subalgebras of $\mathfrak{so}(3,1)$ are $\{0\}$, $\mathfrak{k}$, and
$\mathfrak{so}(3,1)$.

By Step~1, $\mathfrak{h} \neq \{0\}$. Suppose for contradiction
$\mathfrak{h} = \mathfrak{k}$. Then $\mathfrak{h}$ would be invariant
under $\mathrm{Ad}(B_J)$. For any $v \in \mathbb{R}^3$ not parallel
to $\hat{z}$, direct computation in the standard four-dimensional
representation gives
\[
  \mathrm{Ad}(B_J)\, J(v) \;=\; \cosh\rho\, J(v) \;+\; \sinh\rho\,
  K(\hat{z} \times v),
\]
which has a nonzero $\mathfrak{p}$-component since $\hat{z} \times v
\neq 0$ and $\sinh\rho > 0$. Hence $\mathrm{Ad}(B_J)\, J(v) \notin
\mathfrak{k}$, contradicting $\mathrm{Ad}(B_J)$-invariance. So
$\mathfrak{h} = \mathfrak{so}(3,1)$, and $H^0 = \mathrm{SO}^+(3,1)^0$. Since $H \subseteq
\mathrm{SO}^+(3,1)^0$ and $H \supseteq H^0 = \mathrm{SO}^+(3,1)^0$, we
have $H = \mathrm{SO}^+(3,1)^0$.
\end{proof}

\begin{remark}
The argument depends on $A_5$ being a specific finite subgroup of
$\mathrm{SO}(3)$ acting irreducibly on $\mathbb{R}^3$, on the
existence of an adjacent-vertex pair giving a calculable elliptic
commutator, and on $B_J$ being a non-trivial boost. Replacing $A_5$
with the cubic group $S_4$ (orientation-preserving subgroup of $O_h$)
fails on the irreducibility step: $S_4$ acts on $\mathbb{R}^3$ as an
irreducible three-dimensional representation, but the Niven exclusion
of cubic angles gives candidate cosines that cannot be guaranteed
distinct without case analysis. Replacing $\zeta_5$ with $\zeta_3$,
$\zeta_4$, or $\zeta_8$ removes the natural pairing of an icosahedral
spatial symmetry with a real-rapidity boost coming from a unit norm
relation.
\end{remark}


\section{Lorentz protection}

The density theorem has an immediate functional-analytic corollary.

\begin{lemma}\label{lem:protection}
Let $G$ be a topological group acting continuously on a topological
space $F$. Let $H \leq G$ be a dense subgroup, and let
$S: F \to \mathbb{R}$ be continuous and $H$-invariant. Then $S$ is
$G$-invariant.
\end{lemma}

\begin{proof}
Fix $g \in G$ and $\psi \in F$. Density of $H$ in $G$ gives a net
$(h_i)_{i \in I}$ in $H$ with $h_i \to g$. Joint continuity of the
action gives $h_i \cdot \psi \to g \cdot \psi$ in $F$, and continuity
of $S$ gives
\[
  S(g \cdot \psi) \;=\; \lim_i S(h_i \cdot \psi) \;=\; S(\psi)
\]
by $H$-invariance of $S$ along the net. Limits in $\mathbb{R}$ are
unique, so $S(g \cdot \psi) = S(\psi)$ as required. (For the
application to $G = \mathrm{SO}^+(3,1)^0$, which is a Lie group and
hence first-countable, the same argument with sequences in place of
nets suffices.)
\end{proof}

\begin{corollary}[Lorentz protection]\label{cor:lp}
Let $F$ be a topological space carrying a continuous action of
$\mathrm{SO}^+(3,1)^0$, and let $S: F \to \mathbb{R}$ be a continuous
functional invariant under $\Gamma = \langle A_5, B_J\rangle$. Then
$S$ is $\mathrm{SO}^+(3,1)^0$-invariant.
\end{corollary}

\begin{proof}
By Theorem~\ref{thm:density}, $\Gamma$ is dense in
$\mathrm{SO}^+(3,1)^0$. Apply Lemma~\ref{lem:protection} with
$G = \mathrm{SO}^+(3,1)^0$ and $H = \Gamma$.
\end{proof}

The intended application is to a Wilson effective action
$S_{\mathrm{eff}}$ obtained from a substrate field theory with $\Gamma$
in its symmetry group. The corollary applies under the following
explicit assumptions, which we collect for clarity:

\begin{itemize}
\item[(A1)] The microscopic dynamics, integration measure, regulator,
  and coarse-graining map are all $\Gamma$-equivariant.
\item[(A2)] The effective action $S_{\mathrm{eff}}$ is well-defined as
  a functional on a field configuration space $F$ that carries a
  continuous $\mathrm{SO}^+(3,1)^0$-action.
\item[(A3)] $S_{\mathrm{eff}}$ is continuous on $F$ in some topology
  with respect to which the $\mathrm{SO}^+(3,1)^0$-action is
  continuous (e.g.\ Sobolev $H^k$ for $k$ sufficiently large, or
  Schwartz, or any standard field-theory topology supporting the
  intended observables).
\end{itemize}

These assumptions are satisfied by the standard cases of perturbative
quantum field theory: polynomial Lagrangians of finite degree,
renormalizable interactions, and Wilsonian effective actions with
bounded ultraviolet cutoff. As a concrete sufficient choice, $F$ may
be taken as $H^k_{\mathrm{loc}}(\mathbb{R}^{1,3})$ for some $k$ large
enough to make the action functionals continuous in the field;
$\mathrm{SO}^+(3,1)^0$ acts continuously on this space by pullback,
and finite-degree polynomial actions are continuous on it. Under
(A1)--(A3), Corollary~\ref{cor:lp} forces $S_{\mathrm{eff}}$ to be
$\mathrm{SO}^+(3,1)^0$-invariant. No Lorentz-violating operator can
appear in the effective Lagrangian, at any loop order, for any
contribution continuous in (A3).

The corollary fails when (A3) fails. Two ways this can happen are
worth noting. If continuity fails on a measure-zero subset of $F$
(e.g.\ light-cone singularities of correlation functions), the
$\Gamma$-invariance still extends to $\mathrm{SO}^+(3,1)^0$-invariance
on the dense open complement, and the corollary applies after
restriction. If continuity fails generically, for instance because
the effective action is defined only on a discrete subset of $F$ that
is itself a $\Gamma$-orbit but not closed under the dense subgroup
extension, then the dense-subgroup argument cannot upgrade the
invariance, and the corollary does not apply.

This blocks the Collins--Perez--Sudarsky obstruction at its premise:
the mismatch between the substrate symmetry group and the Lorentz
group, on which the obstruction depends, does not exist for $\Gamma$.
The set of $\Gamma$-invariant continuous operators coincides with the
set of $\mathrm{SO}^+(3,1)^0$-invariant continuous operators, so there
is no Lorentz-violating operator to fine-tune away.

\begin{remark}
The corollary extends to the non-perturbative regime in the following
sense. If the configuration space $F$ admits an integer-valued
topological invariant $W : F \to \mathbb{Z}$ (instanton number,
soliton charge, monopole class), then $W$ is locally constant since
$\mathbb{Z}$ carries the discrete topology, and the level sets
$W^{-1}(n)$ are open and pairwise disjoint. Each $\Lambda \in
\mathrm{SO}^+(3,1)^0$ acts on $F$ as a homeomorphism and therefore
preserves the level sets. Within each sector, $S_{\mathrm{eff}}$ is
continuous in the sense of (A3), and Corollary~\ref{cor:lp} applies
sector-wise. The partition function decomposes as $Z(g) = \sum_n
Z_n(g)$ with each $Z_n$ independently $\mathrm{SO}^+(3,1)^0$-invariant.
\end{remark}

\begin{remark}
Corollary~\ref{cor:lp} requires the field configuration space $F$ to
support a continuous $\mathrm{SO}^+(3,1)^0$-action. This excludes
substrates realized as locally finite point sets in Minkowski space,
since no such set is invariant under a noncompact dense Lorentz
subgroup without producing accumulating orbits. The intended
application is to substrate models in which $F$ is a continuum field
space, with the discrete substrate symmetry $\Gamma$ acting via its
embedding in $\mathrm{SO}^+(3,1)^0$ on continuum fields, rather than
on a literal lattice of spacetime points.
\end{remark}


\section{Numerical appendix: Molien series and spatial anisotropy}

The density theorem of Section~3 is an algebraic statement about
$\mathrm{SO}^+(3,1)^0$; it does not by itself constrain the spatial
anisotropy structure of any specific lattice. However, a separate and
classical Molien-series calculation gives a useful spatial-symmetry
fact about the icosahedral rotation group, which we record here as a
numerical check.

For a finite subgroup $G \subset \mathrm{O}(3)$ acting on
$\mathbb{R}^3$, the Molien series
\[
  M_G(t) \;=\; \frac{1}{|G|} \sum_{g \in G} \frac{1}{\det(1 - tg)}
\]
encodes the Hilbert series of the ring of $G$-invariant polynomials.
For the rotational icosahedral group $G = A_5$, Klein \cite{Klein1884}
gave the closed form
\[
  M_{A_5}(t) \;=\; \frac{1 + t^{15}}{(1 - t^2)(1 - t^6)(1 - t^{10})}.
\]
Reading off coefficients gives the dimension count
\[
  \dim\bigl((S^d \mathbb{R}^3)^{A_5}\bigr) \quad
  \text{for } d = 0, 2, 4, 6, 8, 10, 12, 14, 15, 16
  \quad = \quad 1, 1, 1, 2, 2, 3, 4, 4, 1, 5.
\]
Crucially, $\dim((S^2 \mathbb{R}^3)^{A_5}) = \dim((S^4
\mathbb{R}^3)^{A_5}) = 1$, with unique invariants $|k|^2$ and
$(|k|^2)^2$ respectively. The first non-isotropic $A_5$-invariant
appears at degree six: $\dim((S^6 \mathbb{R}^3)^{A_5}) = 2$, with the
two-dimensional space spanned by $(|k|^2)^3$ and the Klein form $K_6$.
For comparison, the cubic point group $O_h$ has $\dim((S^4
\mathbb{R}^3)^{O_h}) = 2$, so cubic-symmetric couplings have
non-isotropic invariants already at degree four.

This Molien-series fact has a direct numerical consequence on graph
Laplacians. For couplings $e_1, \dots, e_m \in \mathbb{R}^3$ closed
under a finite subgroup $G \subset \mathrm{SO}(3)$, the dispersion
relation
\[
  E(k) \;=\; \sum_i \bigl(1 - \cos(k \cdot e_i)\bigr)
\]
expands as a sum of homogeneous polynomials in the components of $k$,
each lying in the corresponding $G$-invariant subspace of $S^d
\mathbb{R}^3$. The angular spread $\Delta(|k|) = \max_n E(|k| n) -
\min_n E(|k| n)$ over the unit sphere is a power-law function of
$|k|$ at small $|k|$, with exponent equal to the smallest degree at
which a non-isotropic $G$-invariant appears.

We test this numerically. Couplings are taken as the six unit vectors
of the cubic axes (sum normalized to common total weight), and the
twelve unit vectors along the icosahedron vertices, normalized
similarly. The unit sphere is sampled by $5000$ directions on a
Fibonacci spiral. Linear regression of $\log\Delta$ on $\log|k|$ in
the regime $|k| \leq 0.3$ gives:

\begin{table}[ht]
\centering
\begin{tabular}{|l|c|c|}
\hline
coupling family & predicted exponent & fitted exponent \\
\hline
cubic ($O_h$)        & 4 & 3.998 \\
icosahedral ($A_5$)  & 6 & 5.998 \\
\hline
\end{tabular}
\caption{Power-law exponent of the angular spread $\Delta(|k|)$ for
two families of unit-vector couplings in $\mathbb{R}^3$, fitted in
the regime $|k| \leq 0.3$. The exponents agree with the predictions
of the Molien series $M_G(t)$ for $G = O_h$ (degree-four anisotropy)
and $G = A_5$ (degree-six anisotropy) to three significant figures.}
\end{table}

The exponents agree with the Molien-series prediction to three
significant figures.

The boundary between the two statements is sharp. The density theorem
of Section~3 protects Lorentz symmetry in the continuum effective
dynamics, given continuity and exact $\Gamma$-symmetry, via
Corollary~\ref{cor:lp}. The numerical check above is a property of
the finite spatial subgroup $A_5 \subset \mathrm{SO}(3)$ acting on
$\mathbb{R}^3$: it shows that $A_5$-invariant lattice couplings
postpone the first angular anisotropy to degree six in the spatial
momentum, two orders later than cubic-symmetric couplings. This is
not a numerical test of Lorentz density, and would not become one if
the fits were extended to higher precision. The Molien-series
prediction and Theorem~\ref{thm:density} are logically independent;
they share only the algebraic origin of the icosahedral group.


\section{Discussion}

Finitely generated dense subgroups of $\mathrm{SO}^+(3,1)^0$ have a
substantial mathematical literature, of which we make no exhaustive
review. Margulis's arithmeticity and superrigidity theorems
\cite{Margulis91} classify lattices in semisimple Lie groups of higher
rank; Ratner's theorems \cite{Ratner91} describe orbit closures and
invariant measures of unipotent flows on homogeneous spaces;
Breuillard and Gelander \cite{BG03} prove in particular that every
connected Lie group of positive dimension admits a finitely generated
dense free subgroup. The construction of the present paper is
specific rather than generic: it gives an explicit cyclotomic
mechanism with computable elliptic commutator and closed-form
rotation cosine in $\mathbb{Q}(\sqrt{5})$.

The construction of Section~2 gives one finitely-generated
algebraic subgroup of $\mathrm{SO}^+(3,1)^0$ that is dense in the
identity component. The proof is
constructive: an explicit elliptic commutator with closed-form
rotation cosine in $\mathbb{Q}(\sqrt{5})$, plus an irreducibility
argument for the rotational icosahedral representation. The structural
features that the proof exploits are: a non-abelian finite rotation
subgroup acting irreducibly on $\mathbb{R}^3$, a boost whose
commutator with one of its conjugates produces an infinite-order
elliptic Lorentz element, and the bracket structure of
$\mathfrak{so}(3,1)$ which forbids three-dimensional Lie subalgebras
mixing rotations and boosts beyond pure rotations.

We do not claim that an irrational rapidity alone suffices. The
powers of a single boost lie in a discrete one-parameter subgroup of
$\mathrm{SO}^+(3,1)^0$ regardless of rationality. The hypothesis that
makes the present proof work is the production of a specific
infinite-order elliptic element via the conjugate-boost commutator,
together with $A_5$-irreducibility ruling out proper
$\mathrm{Ad}(\Gamma)$-invariant Lie subalgebras. Similar constructions
may work for other finite rotation groups and other boosts whenever
the same combination of conditions holds; we do not attempt a
classification here.

The cyclotomic origin of the rapidity is natural: the unit $J = 1 +
\zeta_5^2$ in $\mathbb{Z}[\zeta_5]$ has Galois-conjugate moduli
$(\varphi^{-1}, \varphi, \varphi, \varphi^{-1})$, providing a
canonical real positive logarithm $\log\varphi$ for the boost rapidity.
Cyclotomic fields of smaller conductor do not produce dense subgroups
via the same construction. $\mathbb{Z}[\zeta_3]$ and $\mathbb{Z}[i] =
\mathbb{Z}[\zeta_4]$ have no real-positive units other than units of
finite order, so no irrational rapidity is available. $\mathbb{Z}[\zeta_8]$
has units $1 + \sqrt{2}$ but the associated finite rotation groups
($A_4$, $S_4$ and their double covers) do not place an icosahedral
adjacency at the calculable angle $\arccos(1/\sqrt{5})$.

The Lorentz protection corollary is independent of the specific
example. Any finite-rank algebraic substrate whose symmetry group is
dense in $\mathrm{SO}^+(3,1)^0$ is automatically immune to the
Collins--Perez--Sudarsky obstruction, under the regularity assumptions
(A1)--(A3). The substantive question for any candidate substrate is
whether its symmetry group is dense in $\mathrm{SO}^+(3,1)^0$. This is
a tractable algebraic question, decidable by methods of the type used
here.

The construction does not address the conformal factor of the
emergent metric or the topological sector structure of the
configuration space; these require separate inputs. Density alone
provides Lorentz invariance, not geometry.

A final caveat on the substrate interpretation. The result should not
be read as asserting that an ordinary locally finite spacetime
lattice can carry $\Gamma$ as a group of literal point automorphisms.
A noncompact dense Lorentz subgroup acting on a discrete subset of
Minkowski space necessarily produces accumulating orbits, which is
incompatible with local finiteness. The construction identifies a
dense algebraic symmetry mechanism that lives on a continuum
configuration space; realizing it in a microscopic substrate requires
an additional model-dependent implementation.

We sketch one natural family of realizations to fix ideas. Take field
configurations valued in finite-rank free $\mathcal{O}_K$-modules
$\mathcal{O}_K^d$, with a Galois-equivariant update rule in the
sense that the dynamics commutes with $\mathrm{Gal}(K / \mathbb{Q})$
acting coordinatewise. Multiplication by a unit
$u \in \mathcal{O}_K^\times$ preserves the lattice and embeds in
$\mathrm{End}(\mathcal{O}_K^d)$ as an integer matrix; the unit
$J = 1 + \zeta_5^2$ realizes the boost $B_J$ in the corresponding
real embedding. The icosahedral $A_5$ acts by an integer
representation on the relevant lattice, and $\Gamma$ generates the
substrate symmetry. The continuum effective action obtained by
coarse-graining such a substrate sits naturally in a Sobolev space of
the type discussed for (A1)--(A3), and the Lorentz protection
corollary applies. We do not develop a specific dynamics here; the
point is only that the algebraic mechanism has plausible carriers and
is not vacuously satisfied. The Lorentz protection corollary applies
whenever any such implementation exists and satisfies (A1)--(A3); the
present paper makes no commitment to a particular microscopic model.


\section*{Acknowledgements}

The Molien-series and Niven verifications, the elliptic-commutator
trace calculation, and the numerical dispersion test were carried out
using the \texttt{sympy}, \texttt{mpmath} and \texttt{numpy}
libraries. Symbolic-algebra and proof-search tools were used in
checking intermediate computations. The author thanks technical
reviewers and proof-checking tools for several rounds of close
reading that identified the invalidity of an earlier
polar-decomposition density argument, proposed the present
Lie-algebra normality route, suggested the explicit elliptic
classification (rank and fixed-plane signature) and the explicit
$A_5$ element on icosahedron vertices, and recommended the cyclotomic
cosine exclusion lemma in place of an opaque appeal to Niven's
theorem. Computations are reproducible from the companion artifact
\cite{ALPCode}.


\begin{thebibliography}{99}

\bibitem{CPSUV04}
J.~Collins, A.~Perez, D.~Sudarsky, L.~Urrutia, H.~Vucetich,
\emph{Lorentz invariance and quantum gravity: an additional
fine-tuning problem?},
Phys.\ Rev.\ Lett.\ \textbf{93} (2004) 191301.
DOI: \texttt{10.1103/PhysRevLett.93.191301}.
arXiv:\texttt{gr-qc/0403053}.

\bibitem{Mattingly05}
D.~Mattingly,
\emph{Modern tests of Lorentz invariance},
Living Rev.\ Relativity \textbf{8} (2005) 5.
DOI: \texttt{10.12942/lrr-2005-5}.

\bibitem{LiberatiReview}
S.~Liberati,
\emph{Tests of Lorentz invariance: a 2013 update},
Class.\ Quant.\ Grav.\ \textbf{30} (2013) 133001.
DOI: \texttt{10.1088/0264-9381/30/13/133001}.

\bibitem{BHS87}
L.~Bombelli, J.~Lee, D.~Meyer, R.~D.~Sorkin,
\emph{Space-time as a causal set},
Phys.\ Rev.\ Lett.\ \textbf{59} (1987) 521.

\bibitem{Klein1884}
F.~Klein,
\emph{Vorlesungen \"uber das Ikosaeder und die Aufl\"osung der
Gleichungen vom f\"unften Grade},
Teubner, Leipzig, 1884.

\bibitem{Niven}
I.~Niven,
\emph{Irrational Numbers},
Carus Math.\ Monograph No.\ 11, MAA, 1956.

\bibitem{Margulis91}
G.~A.~Margulis,
\emph{Discrete Subgroups of Semisimple Lie Groups},
Ergebnisse der Mathematik 17, Springer, 1991.

\bibitem{Ratner91}
M.~Ratner,
\emph{On Raghunathan's measure conjecture},
Ann.\ Math.\ \textbf{134} (1991) 545--607.
DOI: \texttt{10.2307/2944357}.

\bibitem{BG03}
E.~Breuillard, T.~Gelander,
\emph{On dense free subgroups of Lie groups},
J.\ Algebra \textbf{261} (2003) 448--467.
DOI: \texttt{10.1016/S0021-8693(02)00665-1}.

\bibitem{ALPCode}
A.~M.~Thorn,
\emph{algebraic-lorentz-protection: verification scripts}, 2026.\\
DOI: \url{https://doi.org/10.5281/zenodo.20029689}.\\
Source: \url{https://github.com/mathorn1973/algebraic-lorentz-protection}.

\end{thebibliography}

\end{document}
